\documentclass[a4paper,twoside]{article}

\usepackage[T1]{fontenc}

% SCITEPRESS/ICAART template packages
\usepackage{epsfig}
\usepackage{subcaption}
\usepackage{calc}
\usepackage{amssymb}
\usepackage{amstext}
\usepackage{amsmath}
\usepackage{amsthm}
\usepackage{amsfonts}
\usepackage{array}
\usepackage{booktabs}
\usepackage{color}
\usepackage{multicol}
\usepackage{pslatex}
\usepackage{apalike}
\usepackage[ruled,vlined]{algorithm2e}
\usepackage[bottom]{footmisc}
\usepackage[hidelinks]{hyperref}
\usepackage{SCITEPRESS}

\newtheorem{theorem}{Theorem}
\newtheorem{lemma}{Lemma}
\newtheorem{proposition}{Proposition}
\newtheorem{example}{Example}
\theoremstyle{definition}
\newtheorem{definition}{Definition}
\newtheorem{nremark}{Remark}

\renewcommand{\T}{\mathcal{T}}
\renewcommand{\R}{\mathbb{R}}
\newcolumntype{L}[1]{>{\raggedright\arraybackslash}p{#1}}

\sloppy

\title{Fair Admissions with Overlapping Reserve Categories under Score Constraints}

% Double-blind review version. Replace with real author data for camera-ready submission.
\author{\authorname{Anonymous Author(s)}
\affiliation{\sup{1}Anonymous Institution}
\email{anonymous@example.com}
}

% =====================================================================
% [AI 2026-08-18] ICAART 2027 投稿ステータス。詳細は TODO.md 参照。
%
%   締切: Regular Papers 2026-09-30 / Full Paper 12ページ
%   分量: 本文の非空白文字数 約 49,568 / 上限 50,000 字（残り約430字）
%         推定 12.5〜13 ページ。全幅フロート table*2 + algorithm*3 が効いている
%         → 実測では超過見込み。削減候補は本文中の [AI: 削減候補] を参照
%   未検証: この作業環境に latexmk/pdflatex が無く、コンパイルできていない。
%           実ページ数・実文字数は要確認
%   投稿パッケージから Example.tex / Example.pdf を除外すること
%
%   バックアップ: main.tex.before_icaart_structure（構成変更前）
%                 main.tex.before_ai_comments（このコメント追加前）
%                 references.bib.before_prior_work_check
% =====================================================================
\begin{document}

% Previous abstract (251 words, over the SCITEPRESS 200-word ceiling) kept for reference:
% We study admissions with one open category and several reserve categories when students may qualify for more than one reserve category at once. Under one-to-one counting, an admitted student is counted toward at most one category, so a mechanism must decide both who is admitted and which category each overlapping student fills. We also impose a score constraint bounding the gap between the open cutoff and each reserve cutoff. Non-wastefulness must be restated relative to this constraint, while within-category fairness need not be weakened because replacing a reserve admit by a higher-scoring eligible student cannot destroy score-feasibility. A student-proposing score-elevated deferred-acceptance mechanism (SPSE) is feasible, non-wasteful, strategy-proof and free of justified envy, but need not be score-feasible. We isolate the additional role of overlap: routing a multiply eligible student can itself determine whether the cutoff rule survives. We then show that score-feasibility, open-completeness and reserve-completeness may be jointly unattainable, through a Hall-type obstruction on non-singleton category sets, so any mechanism enforcing the cutoff rule must abandon one of them. We study the two resulting responses. Threshold-adjusting deferred acceptance (A-STDA) keeps the announced capacities and gives up reserve-completeness, leaving reserve seats idle; it is score-feasible and fair within categories, and it is exactly the mechanism for which the score-relative restatement of non-wastefulness has content. Preference-aware overlapping de-reservation (PA-ODR) instead keeps every surviving seat filled and gives up fixed capacities; it retains within-category fairness but forfeits strategy-proofness and no justified envy. A prototype on randomized instances prices the two responses against each other.
% [AI] SCITEPRESS 規定 70〜200 語。現在 195 語。加筆時は 200 語を超えないこと。
\abstract{
We study admissions with one open category and several reserve categories when a student may qualify for several at once. Under one-to-one counting an admitted student fills at most one category, so a mechanism must decide both who is admitted and which category each fills. We add a score constraint bounding the gap between the open cutoff and each reserve cutoff. Non-wastefulness must then be restated relative to it, and beside within-category fairness we place across-category efficiency, asking that no two admitted students would rather swap the categories they fill. Score-elevated deferred acceptance is fair, non-wasteful and strategy-proof but need not be score-feasible: routing a multiply eligible student can itself decide whether the cutoff rule survives. We then show that score-feasibility, open-completeness and reserve-completeness may be jointly unattainable, through a Hall-type obstruction on non-singleton category sets, and study the two ways out. Adaptive score-threshold deferred acceptance keeps the announced capacities and stays fair, but its threshold overshoots and admits no fixed point, which costs it non-wastefulness and strategy-proofness alike; preference-aware overlapping de-reservation keeps every surviving seat filled and gives up fixed capacities, strategy-proofness, and fairness to any student already seated. A prototype prices the two against each other.
}

\keywords{Matching, Affirmative Action, Reserve Categories, Overlapping Types, De-Reservation, School Choice}

\onecolumn \maketitle \normalsize \setcounter{footnote}{0} \vfill

\section{\uppercase{Introduction}}
Reserve-based admissions and controlled school choice are studied extensively \cite{Abdu05a,EHYY14a,HYY13a,EcYe15a,AyTu20a,AyTu23a}. In standard formulations each student belongs to one protected category or none; in many real systems they overlap. S\"onmez and Yenmez treat overlapping reserves within the vertical--horizontal system of Indian affirmative action \cite{SoYe22a}; related environments appear in \cite{KHIY17a,AyBo21a,DKT23a,CGH20a} and, under distributional constraints, in \cite{GKK+17a,KaKo24a}. Indian engineering admissions, the setting closest to ours, are described in \cite{BCC+19b}. That \emph{which} slot a student is counted against is itself a design variable is explicit in \cite{KoSo16a}, the cost of getting it wrong documented in \cite{DKPS18a}; the processing order of reserve categories is likewise consequential \cite{PSUY24a,PRS23a}, a sensitivity that reappears in PA-ODR.

Once categories overlap, a mechanism must decide not only who is admitted but how overlapping students are counted, since under one-to-one counting such a student fills at most one reserve slot: the category she is assigned is part of the allocation. We add a score constraint keeping the open and reserve cutoffs close. This is not stylized---in \emph{Ashoka Kumar Thakur vs.\ Union of India} the Supreme Court of India recommended a maximum ten-point gap between the open and OBC cutoffs, analyzed in \cite{AyTu23b}, with cutoff concerns more generally in \cite{AyTu23a,ABB24a,HuZh25a}. Under such a rule a reserve category may be nominally available yet unusable.

These ingredients hit the standard desiderata \cite{EHYY14a,KHIY17a,GKK+17a} asymmetrically. Non-wastefulness must be restated, an empty reserve category being possibly unusable, since \emph{adding} a student to a reserve category can pull its cutoff below the permitted gap. Fairness need not be, since \emph{replacing} a reserve admit by a higher-scoring eligible one only raises that cutoff; what fairness must do instead is reach across categories, because with overlap a strong student eligible for several can protect the cutoff in only one.

\subsubsection{Relation to the closest prior work.}
% [AI] 未解決: SoYe22a との novelty 切り分けが未記載。詳細は TODO.md「残りの査読対策」。
S\"onmez and Yenmez \cite{SoYe22a} take reserve capacities as given and impose no cutoff-gap requirement, the constraint driving every negative result here. Ayg\"un and Turhan \cite{AyTu23b} study exactly the score-elevated policy and the ten-point rule, but without overlapping eligibility, so their mechanism never decides \emph{which} category a student fills; their hierarchical treatment of overlapping horizontal categories \cite{AyTu26a} admits multiple qualifications but is one-to-all with no cutoff-gap rule, so again nothing is routed. What is new is thus not that a score-elevated mechanism can breach a cutoff-gap rule---that is the inconsistency of \cite{AyTu23b} and needs no overlap---but a second, independent cause: eligible for several categories yet counted once, a student has her routing decide whether the rule survives. Proposition~\ref{prop:spse-not-sf} separates this cause, and Proposition~\ref{prop:impossibility} excludes the first by hypothesis. The de-reservation response follows \cite{AyTu23a}.

\subsubsection{Contributions.}

% \textcolor{red}{Contribution needs to say: adapt the properties to this setting; use DA/SPSE and show it does not guarantee score-feasibility; state the impossibility result; make PA-ODR/de-reservation the main constructive contribution; include the small simulation/prototype check.}

\begin{enumerate}
\item We adapt the standard properties to the setting: non-wastefulness becomes score-relative (Def.~\ref{def:nw}), while fairness splits into a within-category priority requirement (Def.~\ref{def:fair}) and \emph{across-category efficiency} (Def.~\ref{def:acrosseff}), which asks only that no two matched students would rather swap. The two are independent, and it is the second that overlap makes necessary. We also fix what capacity change means (Section~\ref{sec:model}).
\item We modify student-proposing deferred acceptance for score-elevated categories (SPSE): feasible, non-wasteful, strategy-proof and envy-free, but not guaranteed score-feasible (Section~\ref{sec:spse}).
\item We prove score-feasibility, open-completeness and reserve-completeness jointly unattainable even when each reserve category is individually fulfillable (Proposition~\ref{prop:impossibility}). This is the pivot: enforcing the cutoff rule costs one of the three, and the last two sections are the two ways of paying.
\item A-STDA keeps the announced capacities and gives up reserve-completeness, restricting proposals to score-admissible pairs and raising its threshold to the cutoff its own outcome induces. It is score-feasible and across-category efficient---without the no-justified-envy property from which SPSE gets that---but wasteful even in the score-relative reading and manipulable, both because that threshold overshoots and admits no fixed point (Propositions~\ref{prop:tada-nw}, \ref{prop:tada-notsp}).
\item PA-ODR gives up fixed capacities instead, subject only to conservation of the total. It enforces score-feasibility and keeps non-wastefulness relative to the final capacities, together with fairness to any \emph{unmatched} claimant; it gives up strategy-proofness and both fairness notions as soon as the claimant already holds a seat, already inside the reserve categories (Propositions~\ref{prop:paodr-fair} and~\ref{prop:paodr-unfair}, Section~\ref{sec:paodr}). We implement its reserve-side choice as maximum-weight bipartite matching and report a prototype across capacity regimes (Section~\ref{sec:experiments}).
\end{enumerate}

\begin{table*}[t]
\caption{The three mechanisms. Feasibility is the only property all three keep outright; each of the rest is abandoned by exactly one or two of them, and the last two rows say at what price. SPSE abandons score-feasibility; A-STDA abandons reserve-completeness (Prop.~\ref{prop:impossibility}) and with it non-wastefulness; PA-ODR abandons the announced capacities, conserving their total (Def.~\ref{def:dereserve}). The two fairness notions are independent of each other (Section~\ref{sec:model}).}
\label{tab:comparison}
\centering
\footnotesize
\setlength{\tabcolsep}{4pt}
\begin{tabular}{|L{0.22\textwidth}|L{0.20\textwidth}|L{0.24\textwidth}|L{0.24\textwidth}|}
\hline
\textbf{Property} & \textbf{SPSE} & \textbf{A-STDA} & \textbf{PA-ODR} \\
\hline
Feasibility & Yes, at $q$ & Yes, at $q$ & Yes, at final $\hat q$ \\
Score-feasibility (Def.~\ref{def:scorefeas}) & No (Prop.~\ref{prop:spse-not-sf}) & Yes (Prop.~\ref{prop:tada-props}) & Yes (Prop.~\ref{prop:paodr-basic}) \\
Non-wastefulness (Def.~\ref{def:nw}) & Yes (Prop.~\ref{prop:spse-basic}) & No (Prop.~\ref{prop:tada-nw}) & Yes, w.r.t.\ $\hat q$ \\
Within-category fairness (Def.~\ref{def:fair}) & Yes (Prop.~\ref{prop:spse-nje}) & Yes (Prop.~\ref{prop:tada-props}) & No; unmatched claimants only (Props.~\ref{prop:paodr-fair}, \ref{prop:paodr-unfair}) \\
Across-category efficiency (Def.~\ref{def:acrosseff}) & Yes (Prop.~\ref{prop:spse-nje}) & Yes (Prop.~\ref{prop:tada-props}) & No; yes inside $T$ (Sec.~\ref{sec:paodr-limits}) \\
\hline
Strategy-proofness & Yes (Prop.~\ref{prop:spse-sp}) & No (Prop.~\ref{prop:tada-notsp}) & No (Prop.~\ref{prop:paodr-notsp}) \\
De-reservation & None & None & Yes, and not minimal \\
% Uses category preferences & Through proposals & Through lexicographic tie-breaking \\
% Three-column version of the row above: Uses category preferences & Through proposals & Through proposals, inside $A_i(e)$ & Through lexicographic tie-breaking \\
\hline
\end{tabular}
\end{table*}

\section{\uppercase{Model}}
\label{sec:model}

\subsection{Primitives, Counting, Scores and Cutoffs}

We consider admissions to a single school. Let $I$ be a finite set of students, $T$ a finite set of reserve categories, $o$ the open category, and $V=T\cup\{o\}$. Each $c\in V$ has capacity $q^c\in\mathbb{N}$, with total intake $Q:=\sum_{c\in V}q^c$; each $i\in I$ has a raw score $\sigma(i)\in\R_+$, distinct across students, and a reserve-eligibility set $\T(i)\subseteq T$, with $i$ \emph{overlapping} when $|\T(i)|\ge2$. All reserve categories share a common bonus $k\ge0$, and $\delta\ge0$ is the allowed cutoff gap. The admissible set of $i$ is $V_i:=\{o\}\cup\T(i)$, over which she has a strict preference $\succ_i$, with $r_i(c)$ the rank of $c$ ($1$ = best): with a single school, students rank admission routes rather than schools.

We adopt \emph{one-to-one counting}: a matched student counts toward exactly one category chosen by the mechanism, not toward every one she is eligible for. This is the convention of \cite{KHIY17a,AyBo21a,GKK+17a}: a student qualifying under both a socioeconomic and a geographic criterion enters by either route but cannot fill both quotas.

\begin{definition}[Matching and feasibility]
A matching is a tuple $\mu=(\mu_c)_{c\in V}$ with $\mu_c\subseteq I$; write $M(\mu):=\bigcup_{c\in V}\mu_c$ for the matched students. It is \emph{feasible} if (i) $|\mu_c|\le q^c$ for every $c\in V$, (ii) $\mu_c\cap\mu_{c'}=\emptyset$ for distinct $c,c'$, (iii) $c\in V_i$ for every $i\in\mu_c$, and (iv) $|M(\mu)|\le Q$.
\end{definition}

Condition (ii) is one-to-one counting and (iv) follows from it; $Q$ is recorded separately because it, not its split, is what de-reservation holds fixed (Definition~\ref{def:dereserve}). The open category ranks by raw score, reserve categories by score plus the common bonus: for $c\in V_i$ the \emph{adjusted score} is $\sigma_c(i):=\sigma(i)$ if $c=o$ and $\sigma(i)+k$ if $c=t\in\T(i)$. Reserve categories thus share one ranking, differing only through eligibility, capacities and preferences; $k=0$ is covered throughout.

\begin{definition}[De-reservation and conservation of the total]
\label{def:dereserve}
An \emph{effective capacity vector} is $\hat q=(\hat q^c)_{c\in V}$ with $\hat q^t\le q^t$ for every $t\in T$ and $\sum_{c\in V}\hat q^c=\sum_{c\in V}q^c$, equivalently $\hat q^o=q^o+\sum_{t\in T}(q^t-\hat q^t)$. A mechanism \emph{de-reserves} when it operates at some $\hat q\neq q$; feasibility, cutoffs and non-wastefulness are read against the vector it operates at, which is $\hat q=q$ for Sections~\ref{sec:spse} and~\ref{sec:tada}.
\end{definition}

The only discipline we impose is conservation of the total: seats are relabelled from reserve to open, never created or destroyed. The enrolment target is an external commitment and the split an internal instrument, so judging a mechanism at $\hat q$ costs the quota and not the seat---the entering class is as large, only the route in has moved. The inequality $\hat q^t\le q^t$ makes this one-directional; Section~\ref{sec:experiments} shows the cost in quota is real nonetheless.

\begin{definition}[Cutoffs and score-feasibility]
\label{def:scorefeas}
For a matching $\mu$ let $e_o(\mu):=\min_{i\in\mu_o}\sigma(i)$, with $e_o(\mu)=-\infty$ if $\mu_o=\emptyset$, and for $t\in T$ let $e_t^+(\mu):=\min_{i\in\mu_t}\sigma_t(i)$, with $e_t^+(\mu)=+\infty$ if $\mu_t=\emptyset$. A feasible $\mu$ is \emph{score-feasible} if $e_o(\mu)-e_t^+(\mu)\le\delta$ for every $t\in T$.
\end{definition}

We read $e_o(\mu)-e_t^+(\mu)$ as $-\infty$ whenever $\mu_o$ or $\mu_t$ is empty, so the condition is vacuous there. No reserve cutoff may fall too far below the open one on the adjusted scale: the bonus is institutionally legitimate and $\delta$ bounds only the gap left after it, so the two parameters do different work. In raw scores they compound, a reserve admit falling up to $\delta+k$ below the open cutoff; for $k=0$, $\delta$ alone bounds the raw gap, the ten-point rule of \cite{AyTu23b}. Score-feasibility is an ex post requirement on realized cutoffs, not implied by the ranking rule, which is why Sections~\ref{sec:tada} and~\ref{sec:paodr} enforce it directly.

\subsection{Desiderata}

Wastefulness is usually defined by an empty category an unmatched student could occupy; here she can occupy it only when score-feasibility is preserved.

\begin{definition}[Non-wastefulness]
\label{def:nw}
A feasible matching $\mu$ is \emph{non-wasteful} if there are no unmatched $i\in I\setminus M(\mu)$ and $c\in V_i$ such that adding $i$ to $c$ yields a matching that is both feasible and score-feasible.
\end{definition}

So a reserve category may be empty yet unusable, and only a student nobody has seated may complain of it. The definition applies to feasible, not only score-feasible, matchings, since SPSE is non-wasteful without being score-feasible; and it requires the \emph{enlarged} matching to be score-feasible, not score-feasibility to be preserved, which would make an already infeasible matching vacuously non-wasteful. The clause has content because adding a student to a reserve category can lower its cutoff; \emph{replacing} an admit by a higher-scoring eligible one cannot, which is why the two fairness notions carry no such clause.

Fairness comes in two kinds: one compares scores inside a single category, the other compares preferences across two. Write $\mu(i)$ for the category holding $i$, and read $c\succ_i\mu(i)$ as true when $i$ is unmatched.

\begin{definition}[Within-category fairness]
\label{def:fair}
A feasible matching $\mu$ is \emph{fair within categories} if there are no $t\in T$ and students $i,j$ with (i) $t\in\T(i)\cap\T(j)$, (ii) $\sigma_t(i)>\sigma_t(j)$, (iii) $j\in\mu_t$, and (iv) $t\succ_i\mu(i)$.
\end{definition}

Clause (iv) covers an unmatched claimant and a matched one alike, so a student holding a seat may point at a category she ranks higher. The distinction is not idle: PA-ODR fails Definition~\ref{def:fair} and keeps exactly its restriction to unmatched claimants (Propositions~\ref{prop:paodr-fair} and~\ref{prop:paodr-unfair}). Preferences alone, with no score comparison, support the second kind.

\begin{definition}[Across-category efficiency]
\label{def:acrosseff}
A feasible matching $\mu$ is \emph{across-category efficient} if there are no $i,j\in M(\mu)$ with (i) $\mu(j)\in V_i$ and $\mu(i)\in V_j$, (ii) $\mu(j)\succ_i\mu(i)$ and $\mu(i)\succ_j\mu(j)$, and (iii) the matching that exchanges $i$ and $j$ feasible and score-feasible.
\end{definition}

This is an exchange requirement rather than a priority one, and it is what makes the routing of a strong overlapping student visible to the two students it inconveniences; it also reaches across the open category, about which Definition~\ref{def:fair} says nothing. Neither notion implies the other. The instance of Proposition~\ref{prop:paodr-unfair} is across-category efficient and not fair within categories; conversely, with $T=\{t\}$, $q^o=q^t=1$, $k=0$ and $\delta=10$, seating $i(80,\{t\})$ at $o$ and $j(90,\{t\})$ at $t$ where $t\succ_i o$ and $o\succ_j t$ is fair within categories, $j$ outscoring $i$ at $t$, while the two block by exchange, the swap leaving $e_o-e_t^+=10\le\delta$.

A circularity underlies everything: feasible reserve assignments depend on the open cutoff, which depends on who is left after reserve assignment. Score-feasibility is thus no static eligibility condition, and greedy quota-filling is unreliable.

\section{\uppercase{Student-Proposing Score-Elevated Mechanism}}
\label{sec:spse}

% \textcolor{red}{Section summary: DA/SPSE satisfies the usual DA properties, but it does not guarantee score-feasibility. This point should be visible before the propositions.}

Students apply to eligible categories and each keeps the highest adjusted-score applicants. SPSE inherits the usual DA advantages---feasibility, non-wastefulness, fairness, strategy-proofness---but may breach score-feasibility.

\begin{algorithm}[t]
\footnotesize
\caption{Student-Proposing Score-Elevated Mechanism (SPSE)}
\label{alg:spse}
\KwIn{$I$, $V=T\cup\{o\}$, capacities $(q^c)_{c\in V}$, preferences $(\succ_i)_{i\in I}$, scores $\sigma$, eligibility $\T(\cdot)$, bonus $k$}
\KwOut{Matching $\mu=(\mu_c)_{c\in V}$}
$\mu_c\gets\emptyset$ for all $c\in V$; mark every student unmatched\;
Set $V_i:=\{o\}\cup\T(i)$ and compute $\sigma_c(i)$ for every $i$ and every $c\in V_i$\;
\While{some unmatched $i$ still has an unapplied category in $V_i$}{
    $i$ applies to her most preferred unapplied $c\in V_i$\;
    $c$ considers $\mu_c\cup\{i\}$ and keeps at most $q^c$ students with highest $\sigma_c$\;
    Every rejected student becomes unmatched\;
}
\Return{$\mu$}\;
\end{algorithm}

SPSE preserves one-to-one counting, a student being held by at most one category at a time, and differs from ordinary student-proposing deferred acceptance only in proposing to categories and in the score elevation. Each student applies at most once per category, so the loop runs at most $|I|(|T|+1)$ times.

\begin{proposition}[Feasibility and non-wastefulness]
\label{prop:spse-basic}
The output of SPSE is feasible and non-wasteful in the sense of Definition~\ref{def:nw}.
\end{proposition}

\begin{proof}
Each category keeps at most its capacity, no student is held twice and $i$ applies only within $V_i$: feasibility. For non-wastefulness, suppose unmatched $i$ could be added to $c\in V_i$; that needs $|\mu_c|<q^c$. Being unmatched at termination, $i$ applied to every category of $V_i$ and was rejected. A category rejects only when full and thereafter only replaces a held student by a strictly better one, so $|\mu_c|=q^c$, a contradiction.
\end{proof}

\begin{proposition}[Fairness of SPSE]
\label{prop:spse-nje}
The output of SPSE is fair within categories and across-category efficient.
\end{proposition}

\begin{proof}
Let $\mu$ be the output. We first show that whenever $c\in V_i$ satisfies $c\succ_i\mu(i)$, every $j\in\mu_c$ has $\sigma_c(j)>\sigma_c(i)$. Proposing in $\succ_i$-order and held only until rejected, $i$ can end unmatched, or below $c$ in her own ranking, only after proposing to $c$ and being rejected. At that rejection $c$ held $q^c$ students and dropped the lowest-$\sigma_c$ member of $\mu_c\cup\{i\}$, so every retained student beat $\sigma_c(i)$; thereafter $c$ stays full and its holdings only improve in $\sigma_c$.

Within-category fairness is this claim at $c=t\in T$, a blocking triple requiring some $j\in\mu_t$ with $\sigma_t(i)>\sigma_t(j)$. For across-category efficiency let $i,j$ block by exchange; the claim at $(i,\mu(j))$ gives $\sigma_{\mu(j)}(j)>\sigma_{\mu(j)}(i)$ and at $(j,\mu(i))$ gives $\sigma_{\mu(i)}(i)>\sigma_{\mu(i)}(j)$. Inside one category the bonus is added to both students or to neither, so each inequality reduces to the same comparison of raw scores and the two contradict each other.
\end{proof}

Raw scores and eligibility are verifiable, so the only object a student can misreport is $\succ_i$.

\begin{proposition}[Strategy-proofness]
\label{prop:spse-sp}
No student can obtain a strictly $\succ_i$-better assignment under SPSE by reporting an ordering over $V_i$ other than $\succ_i$, holding others' reports fixed.
\end{proposition}

\begin{proof}
Treat each $c\in V$ as a college of capacity $q^c$ whose priority over the students with $c\in V_i$ is the strict ranking induced by $\sigma_c$, well defined as raw scores are distinct and responsive since $c$ retains the $q^c$ highest available. Each $i$ ranks all of $V_i$ acceptable, so this is the college-admissions model of Gale and Shapley \cite{GaSh62a} with Algorithm~\ref{alg:spse} its student-proposing deferred acceptance; strategy-proofness for the proposing side is Dubins and Freedman \cite{DuFr81a} and Roth \cite{Roth82a}.
\end{proof}

\begin{proposition}[SPSE does not guarantee score-feasibility]
\label{prop:spse-not-sf}
The output of SPSE need not be score-feasible, even when a matching that is simultaneously score-feasible, non-wasteful and fair within categories exists.
\end{proposition}

\begin{proof}
Let $T=\{t_1,t_2\}$, $q^o=q^{t_1}=q^{t_2}=1$, $k=5$, $\delta=10$, and take $i_1(100,\emptyset)$, $i_2(95,\{t_1,t_2\})$, $i_3(82,\{t_2\})$ with $t_1\succ_{i_2}t_2\succ_{i_2}o$ and $t_2\succ_{i_3}o$, so $\sigma_{t_1}(i_2)=\sigma_{t_2}(i_2)=100$ and $\sigma_{t_2}(i_3)=87$. In round~1 each category receives one applicant and holds her, so SPSE terminates with $\mu_o=\{i_1\}$, $\mu_{t_1}=\{i_2\}$, $\mu_{t_2}=\{i_3\}$ and $e_o-e_{t_2}^+=13>\delta$.

Consider instead $\nu_o=\{i_1\}$, $\nu_{t_2}=\{i_2\}$, $\nu_{t_1}=\emptyset$, where $e_o(\nu)=e_{t_2}^+(\nu)=100$ and the empty $t_1$ is vacuous, so $\nu$ is score-feasible. It is non-wasteful, every category in $V_{i_3}=\{t_2,o\}$ being at capacity, and fair within categories, the sole member of $t_2$ satisfying $\sigma_{t_2}(i_2)=100>87$.
\end{proof}

We do not claim $\nu$ dominates: SPSE matches three students and $\nu$ two, and $\nu$ leaves $t_1$ vacant though $i_2$ ranks it first. The point is that the other properties do not force score-feasibility, an independent requirement bought here at the price of a seat.

The instance also isolates the cause. A score-elevated policy can breach a cutoff-gap rule with no overlap at all, the sole eligible candidate being too weak; that is the inconsistency of \cite{AyTu23b}, and with $T=\{t_1\}$, $i_1(100,\emptyset)$ and $i_2(80,\{t_1\})$ preferring $t_1$, SPSE already returns a gap of $15>\delta$. Here $i_2$ has the same adjusted score in either category and alone would satisfy the rule in whichever she filled, so it is the routing of her single count, not the pool, that breaks it. Under single eligibility there is nothing to route; Proposition~\ref{prop:impossibility} excludes the first cause by hypothesis.

\section{\uppercase{An Impossibility Driven by Overlap}}
\label{sec:impossibility}

An impossibility is easy to obtain for the uninteresting reason just described: in the two-student instance above, open-completeness forces $e_o=100$, nobody clears the bar, and reserve-completeness already fails. We build the negation of that reason into the hypothesis.

\begin{definition}[Individual completability]
\label{def:indcomp}
An instance is \emph{individually completable} if for every $t\in T$ there is a feasible matching that is score-feasible, satisfies open-completeness, and has $|\mu_t|=q^t$, where open-completeness requires $\mu_o$ to consist of the $q^o$ highest raw-score students in $I$.
\end{definition}

So every announced quota is separately honourable in full, score-feasibly and without disturbing the open category; the impossibility is that the promises cannot be kept together.

\begin{proposition}[Impossibility under full capacity utilization]
\label{prop:impossibility}
There is an individually completable instance (Definition~\ref{def:indcomp}) admitting no feasible matching that simultaneously satisfies score-feasibility, \emph{open-completeness} and \emph{reserve-completeness}, where reserve-completeness requires $|\mu_t|=q^t$ for every $t$ for which at least $q^t$ students are eligible.
\end{proposition}

\begin{proof}
Let $T=\{t_1,t_2\}$, $q^o=q^{t_1}=q^{t_2}=1$, $k=5$, $\delta=10$, and take $i_1(100,\emptyset)$, $i_2(95,\{t_1,t_2\})$, $i_3(80,\{t_1\})$, $i_4(78,\{t_2\})$. Open-completeness forces $\mu_o=\{i_1\}$, so $e_o(\mu)=100$ and score-feasibility requires $\sigma_t(i)\ge90$ for every reserve assignee. Since $\sigma_{t_1}(i_2)=\sigma_{t_2}(i_2)=100\ge90$ while $\sigma_{t_1}(i_3)=85<90$ and $\sigma_{t_2}(i_4)=83<90$, student $i_2$ is the only one placeable in any reserve category.

It is individually completable: $\mu_o=\{i_1\}$, $\mu_{t_1}=\{i_2\}$, $\mu_{t_2}=\emptyset$ is open-complete with $|\mu_{t_1}|=q^{t_1}$ and a gap of $0$, the empty $t_2$ being vacuous, and exchanging $t_1$ and $t_2$ handles $t_2$. Yet both quotas cannot be met at once: the eligible sets are $\{i_2,i_3\}$ and $\{i_2,i_4\}$, so reserve-completeness demands two distinct reserve placements while one-to-one counting lets $i_2$ count toward one. Some reserve category must then admit $i_3$ or $i_4$, breaching score-feasibility.
\end{proof}

Both axioms are institutional commitments rather than normative optima: the open category is genuinely open, and an announced quota is honoured whenever the eligible population suffices---how policy documents write quotas, category by category, precisely because they are drafted without reference to overlap. Weakening reserve-completeness to what is \emph{jointly} fillable would restore consistency but concede the point, that overlap silently reduces what a stated quota delivers. The obstruction is worth naming: in the feasibility graph of Section~\ref{sec:paodr} at $e=100$ the two categories share the neighbourhood $\{i_2\}$, so the capacitated Hall condition $|N(S)|\ge\sum_{t\in S}q^t$ fails at $S=T$ while holding on singletons---which is individual completability. That gap is exactly overlap, and it is why PA-ODR is built on a bipartite matching: an obstruction on sets of categories is invisible to a rule processing them one at a time.

\section{\uppercase{Adaptive Score-Threshold Deferred Acceptance}}
\label{sec:tada}

Read in realized cutoffs, Proposition~\ref{prop:impossibility} leaves three ways out: abandon open-completeness, reserve-completeness, or the announced capacities. This section takes the second and the next the third. A-STDA keeps $q$ untouched and declines to fill a reserve seat it cannot fill legally, so the seat stays empty and on the books.

For a threshold $\tau\in\R\cup\{-\infty\}$, let the \emph{admissible application set} of $i$ be
\[
A_i(\tau):=\{o\}\cup\{t\in\T(i):\sigma(i)+k\ge\tau-\delta\},
\]
which depends on scores and eligibility but not on preferences; under a common bonus it contains all of $\T(i)$ or none of it. Student-proposing deferred acceptance with $A_i(\tau)$ in place of $V_i$ and capacities $q$ yields $\mathrm{DA}(\tau)$. \emph{Adaptive score-threshold deferred acceptance} (A-STDA), Algorithm~\ref{alg:tada}, raises $\tau$ to the open cutoff its own outcome induces and stops once that outcome no longer raises it, so each round admits exactly those students the previous round's cutoff permits. Adapting an eligibility rule to the cutoff it produces is the device behind dynamic reserves \cite{AyTu20a}; here it is applied to the score-admissible sets, every round being student-proposing deferred acceptance \cite{GaSh62a,Roth08a}.

\begin{algorithm}[t]
\footnotesize
\caption{Adaptive Score-Threshold Deferred Acceptance (A-STDA)}
\label{alg:tada}
\KwIn{$I$, $V=T\cup\{o\}$, capacities $(q^c)_{c\in V}$, preferences $(\succ_i)_{i\in I}$, scores $\sigma$, eligibility $\T(\cdot)$, bonus $k$, gap $\delta$}
\KwOut{Matching $\mu=(\mu_c)_{c\in V}$}
$\tau\gets-\infty$\;
\While{true}{
    $A_i(\tau)\gets\{o\}\cup\{t\in\T(i):\sigma(i)+k\ge\tau-\delta\}$ for every $i$\;
    $\mu\gets$ student-proposing DA on $(A_i(\tau))_{i\in I}$ with capacities $q$\;
    \lIf{$e_o(\mu)\le\tau$}{\Return{$\mu$}}
    $\tau\gets e_o(\mu)$\;
}
\end{algorithm}

\begin{proposition}[Properties of A-STDA]
\label{prop:tada-props}
A-STDA terminates after at most $|I|+1$ deferred-acceptance runs. Its output is feasible with respect to the original capacities $q$, score-feasible in the realized cutoffs of Definition~\ref{def:scorefeas}, fair within categories and across-category efficient (Definitions~\ref{def:fair} and~\ref{def:acrosseff}). Moreover no unmatched student can be added to the open category, so any waste it leaves is at a reserve category.
\end{proposition}

\begin{proof}
\emph{Termination.} The threshold changes only by rising to $e_o(\mu)$, a raw score, so it never repeats a value and after at most $|I|$ updates the test $e_o(\mu)\le\tau$ must succeed. Write $\tau^\ast$ for the terminal threshold, so $\mu=\mathrm{DA}(\tau^\ast)$ and $e_o(\mu)\le\tau^\ast$.

\emph{Feasibility} is immediate: deferred acceptance holds at most $q^c$ students at $c$ and none at two categories, and proposals run inside $A_i(\tau^\ast)\subseteq V_i$. \emph{Score-feasibility}: every $i\in\mu_t$ proposed to $t$, so $\sigma_t(i)\ge\tau^\ast-\delta\ge e_o(\mu)-\delta$, and empty categories are vacuous by Definition~\ref{def:scorefeas}.

\emph{Fairness.} We claim that whenever $c\in V_i$ satisfies $c\succ_i\mu(i)$, every $j\in\mu_c$ has $\sigma_c(j)>\sigma_c(i)$. If $c\in A_i(\tau^\ast)$, then $i$ proposes in $\succ_i$-order inside $A_i(\tau^\ast)$ and ends worse off than $c$, so she proposed to $c$ and was rejected; a category rejects only when full and afterwards only improves. If instead $c\notin A_i(\tau^\ast)$, then $c$ is a reserve category $t$ with $\sigma(i)+k<\tau^\ast-\delta$, since $o\in A_i(\tau^\ast)$ always; but $j$ proposed to $t$, so $\sigma(j)+k\ge\tau^\ast-\delta>\sigma(i)+k$. Within-category fairness is the claim at $c=t\in T$. For across-category efficiency, let $i,j$ block by exchange; the claim at $(i,\mu(j))$ and at $(j,\mu(i))$ gives $\sigma_{\mu(j)}(j)>\sigma_{\mu(j)}(i)$ and $\sigma_{\mu(i)}(i)>\sigma_{\mu(i)}(j)$, which under the common bonus contradict each other.

\emph{The open category.} $o\in A_i(\tau^\ast)$ for every $i$, so an unmatched $i$ proposed to $o$ and was rejected, giving $|\mu_o|=q^o$.
\end{proof}

The case analysis runs at $\tau^\ast$ itself, needing no property of it beyond the terminal test---which matters, as Proposition~\ref{prop:tada-nw} shows none is available. The second case is why fairness survives the restriction of proposals: a student the threshold silences is one every reserve admit already outscores. The rest of the deferred-acceptance package A-STDA does not inherit.

Away from the open category non-wastefulness fails. At a reserve category $t\notin A_i(\tau^\ast)$ the student was barred from proposing, and adding her gives $e_t^+\le\sigma(i)+k<\tau^\ast-\delta$, so the enlarged matching is score-feasible only if $e_o(\mu)<\tau^\ast$: Definition~\ref{def:nw} holds outright whenever the loop stops at $e_o(\mu)=\tau^\ast$. But the terminal test also accepts $e_o(\mu)<\tau^\ast$, and then a whole band of students is stranded.

\begin{proposition}[A-STDA is wasteful, and the threshold has no fixed point]
\label{prop:tada-nw}
The output of A-STDA need not satisfy Definition~\ref{def:nw}, already with $|T|=1$. In the same instance no threshold satisfies $\tau=e_o(\mathrm{DA}(\tau))$.
\end{proposition}

\begin{proof}
Let $T=\{r,s\}$, $q^o=2$, $q^r=q^s=1$, $k=0$, $\delta=5$, with $c(95,\emptyset)$, $x(88,\{r\})$ ranking $r$ first and $y(86,\{s\})$ ranking $s$ first. At $\tau=-\infty$ deferred acceptance gives $\mu_o=\{c\}$, $\mu_r=\{x\}$, $\mu_s=\{y\}$, a gap of $95-86=9>\delta$, so $\tau\gets95$. At $\tau=95$ the bar is $90$, which neither $88$ nor $86$ clears, so both propose only to $o$, which keeps $c$ and $x$ and rejects $y$: $\mu_o=\{c,x\}$, $\mu_r=\mu_s=\emptyset$ and $e_o(\mu)=88\le\tau$, so A-STDA stops. But $y$ is unmatched, $s$ is empty, and adding her leaves $e_o(\mu)-\sigma_s(y)=88-86=2\le\delta$, an enlargement that is feasible and score-feasible.

For the second claim, $\mathrm{DA}(\tau)$ passes through three regimes as $\tau$ rises: for $\tau\le91$ both clear the bar and $e_o=95$; for $91<\tau\le93$ only $x$ does, $y$ falls to $o$ and $e_o=86$; for $\tau>93$ neither does and $e_o=88$. None of $95$, $86$, $88$ lies in the regime that produces it, so $e_o(\mathrm{DA}(\tau))\neq\tau$ for every $\tau$.
\end{proof}

The threshold overshoots. The bar $\tau^\ast-\delta$ is set from an open cutoff that the students it excludes then push back down, by falling into the open category, and everyone in the resulting band $[e_o(\mu)-\delta,\ \tau^\ast-\delta)$ is left barred from a reserve category that could legally hold her. Iterating $\tau\gets e_o(\mathrm{DA}(\tau))$ to a fixed point is the obvious repair and the second claim rules it out: no fixed point need exist, and the iteration cycles. What remains is to weaken the guarantee to instances stopping at $e_o(\mu)=\tau^\ast$, or to add a pass that seats the band by $\sigma$ after the loop, which cannot break score-feasibility as it only lowers reserve cutoffs. This is also why Proposition~\ref{prop:tada-props} was proved from the terminal test alone. PA-ODR's cutoff is the realized one, so the question does not arise there.

% [AI] Glit26a の incentive 分析が示唆を与えないか要確認（TODO.md 参照）。
Nor is A-STDA strategy-proof, and it fails for the same overshoot. At a \emph{fixed} threshold the mechanism is deferred acceptance on the sets $A_i(\tau)$, which do not depend on preferences, hence strategy-proof by Proposition~\ref{prop:spse-sp}; but $\tau^\ast$ is the open cutoff of a deferred-acceptance outcome, so a student can move it.

\begin{proposition}[A-STDA is not strategy-proof]
\label{prop:tada-notsp}
There is an instance in which a student unmatched under truthful reporting is matched under A-STDA by demoting her true first choice.
\end{proposition}

\begin{proof}
In the instance of Proposition~\ref{prop:tada-nw} let $y$ report $o\succ s$. At $\tau=-\infty$ she proposes to $o$, which holds $c$ and $y$, while $r$ holds $x$; then $e_o=86$ and $e_r^+=88$, the gap is negative, and $\tau\gets86$. At $\tau=86$ the bar is $81$, every proposal is admissible, the same outcome recurs and $e_o=86\le\tau$ stops the loop with $y$ at $o$. Truthfully she is unmatched, so the misreport makes her strictly better off.
\end{proof}

The mechanism is manipulable exactly where it is wasteful: naming the reserve category she truly wants drives $\tau^\ast$ up to the bar that excludes her from it, while reporting $o$ first keeps the threshold low enough to seat her. Both failures are the same overshoot, so what A-STDA pays for keeping $q$ fixed is not confined to coverage.

\section{\uppercase{Preference-Aware Overlapping De-Reservation}}
\label{sec:paodr}

% \textcolor{red}{Add a brief summary of what PA-ODR does and which properties it satisfies/fails.}

PA-ODR fixes the open cutoff, assigns reserves only along score-feasible pairs, and de-reserves unfulfillable capacity to the open category. It is score-feasible but changes capacities and loses the DA-style incentive and envy guarantees. Modifying announced capacities to restore existence is the near-feasibility paradigm, which is known to invite exactly the incentive loss of Section~\ref{sec:paodr-limits} \cite{Glit26a}.

\subsection{Feasibility Graph and Lexicographic Matching}

For a fixed open cutoff $e$, let $G(e)$ be the bipartite graph with left side the students not in $\mu_o$, right side $T$, and an edge $(i,t)$ iff $t\in\T(i)$ and $\sigma(i)+k\ge e-\delta$; under a common bonus a student is score-admissible at $e$ for all her categories or none. Write $P(e)$ for the edge set and fix an injective $p:P(e)\to\{1,\dots,|P(e)|\}$ from linear orders on $I$ and $T$. Among capacitated matchings on $G(e)$, \textsc{PrefLexMaxMatching} is the one selected by, in order: (1)~maximize the number of matched pairs; (2)~lexicographically maximize the weakly decreasingly sorted vector of matched elevated scores $\sigma_t(i)=\sigma(i)+k$; (3)~lexicographically minimize the weakly decreasingly sorted vector of preference ranks $r_i(t)$; (4)~minimize $\sum_{(i,t)\in\mu}2^{p(i,t)}$.

Criterion 3 compares matchings by worst preference rank first, so it is egalitarian. Criterion 4 is a needed tie-breaker: two students eligible for the same two unit-capacity categories and ranking them alike give identical cardinality, score and rank vectors, while distinct edge sets give distinct sums of powers of two. Criteria 1--2 pin down \emph{who} is admitted, leaving 3--4 to decide \emph{where}.

\begin{lemma}[Criteria 1 and 2 determine the admitted set]
\label{lem:matroid}
Fix $e$ and right-side capacities $(\hat q^t)_{t\in T}$. All capacitated matchings on $G(e)$ satisfying criteria 1 and 2 match exactly the same set of students; criteria 3 and 4 only choose its assignment to categories.
\end{lemma}

\begin{proof}
Replace each $t$ by $\hat q^t$ parallel copies, turning capacitated matchings into ordinary ones. The matchable student sets are the independent sets of the induced transversal matroid, and criterion 1 restricts to its bases. Weight $i$ by $\sigma(i)+k$, pairwise distinct as raw scores are. Criterion 2 selects the bases whose decreasingly sorted weight vector is lexicographically maximal, which for distinct weights is the output of the greedy algorithm scanning students in decreasing weight; that algorithm is deterministic, so the basis is unique.
\end{proof}

So preferences matter only after size and score are optimized. The matroid structure underlies \cite{Flei03a} and the substitutable choice functions of \cite{GKK+17a,KoSo16a}; specific here is a ground set defined by an endogenous cutoff. \textsc{PrefLexMaxMatching} runs in polynomial time: weighting each admissible pair so the four criteria occupy decreasing orders of magnitude---possible with $O(|I||T|)$-bit weights---reduces it to one maximum-weight capacitated bipartite matching. By Lemma~\ref{lem:matroid} an implementation may instead take the admitted set by greedy and solve a smaller min-cost assignment, as our prototype does.

\begin{algorithm}[t]
\footnotesize
\caption{Preference-Aware Overlapping De-Reservation (PA-ODR)}
\label{alg:paodr}
\KwIn{$I$, scores $\sigma$, eligibility $\T(\cdot)$, preferences $(\succ_i)_{i\in I}$, capacities $q^o,(q^t)_{t\in T}$, bonus $k$, gap $\delta$, a fixed linear order $\prec$ on $T$}
\KwOut{Matching $\mu=(\mu_o,(\mu_t)_{t\in T})$}
$\hat q^o\gets q^o$ and $\hat q^t\gets q^t$ for all $t\in T$\;
\While{true}{
    Assign the top $\min\{\hat q^o,|I|\}$ raw-score students to $o$; call this set $\mu_o$\;
    $e\gets\min_{i\in\mu_o}\sigma(i)$ if $\mu_o\neq\emptyset$, else $e\gets-\infty$\;
    Build $G(e)$ with edge $(i,t)$ iff $t\in\T(i)$ and $\sigma(i)+k\ge e-\delta$, giving $t$ capacity $\hat q^t$\;
    Compute $(\mu_t)_{t\in T}$ by \textsc{PrefLexMaxMatching}\;
    \eIf{there exists $t$ with $|\mu_t|<\hat q^t$}{
        Let $t^\ast$ be the $\prec$-first such category\;
        $\hat q^{t^\ast}\gets\hat q^{t^\ast}-1$ and $\hat q^o\gets\hat q^o+1$\;
    }{
        \Return{$\mu=(\mu_o,(\mu_t)_{t\in T})$}\;
    }
}
\end{algorithm}

Algorithm~\ref{alg:paodr} builds PA-ODR on \textsc{PrefLexMaxMatching}. It is a fallback mechanism, not a quota-preserving reserve rule: the open category is filled first by raw score, reserve assignment runs only over pairs feasible at the induced cutoff, and capacity a reserve category cannot fill is de-reserved. Merit ordering at $o$ is preserved and preferences only break ties; the order $\prec$ makes the output well defined. PA-ODR is preference-aware only on the reserve side, so a student preferring a reserve category to $o$ may still be placed in $o$, costing it both envy-freeness and strategy-proofness (Section~\ref{sec:paodr-limits}).

Each non-terminating iteration raises $\hat q^o$ by one, so the open sets are nested and the cutoffs $e$ fall weakly. Two opposing forces follow: a lower $e$ only adds edges to $G(e)$, but the same step moves one more student into $\mu_o$. Neither dominates, so a de-reservation can fail to help---the step that frees one category can absorb the only student eligible for another. How many categories PA-ODR retires therefore depends on $\prec$, and it does not minimize de-reservation.

\subsection{Properties of PA-ODR}

With de-reservation in the model (Definition~\ref{def:dereserve}), feasibility and non-wastefulness are judged against the final vector $\hat q$; Section~\ref{sec:experiments} reports how many quotas this retires, so the qualifier is a genuine weakening.

\begin{definition}[Non-wastefulness relative to the final effective capacities]
\label{def:nw-hat}
Let $\hat q$ be the final effective capacity vector produced by PA-ODR. A matching $\mu$ is \emph{non-wasteful relative to $\hat q$} if it satisfies Definition~\ref{def:nw} with feasibility read at $\hat q$.
\end{definition}

\begin{proposition}[Termination, running time, feasibility, score-feasibility, non-wastefulness]
\label{prop:paodr-basic}
PA-ODR terminates; after replacing each $q^c$ by $\min\{q^c,|I|\}$ it runs in time polynomial in the input size. Its output is feasible with respect to the final effective capacities, score-feasible, and non-wasteful relative to those capacities (Definition~\ref{def:nw-hat}).
\end{proposition}

\begin{proof}
\emph{Termination.} A non-terminating iteration moves one seat from a reserve category to $o$, so the nonnegative integer $\sum_t\hat q^t$ falls by one each time and the loop stops after at most $\sum_t q^t$ iterations.

\emph{Running time.} That bound is exponential if capacities are binary-encoded, so we assume $q^c\le|I|$ for every $c\in V$; this is an assumption rather than a normalization, changing how many de-reservation steps the loop takes. Then $\sum_t\hat q^t\le|I||T|$, giving $O(|I||T|)$ iterations, each one sort, one $G(e)$ and one \textsc{PrefLexMaxMatching}.

\emph{Feasibility.} Each iteration gives $o$ the top $\min\{\hat q^o,|I|\}$ students and assigns the rest within $\hat q^t$ along edges of $G(e)$, which join $i$ only to $t\in\T(i)$: no student is assigned twice and no category is over capacity.

\emph{Score-feasibility.} Let $e$ be the final cutoff, so $e_o(\mu)=e$. Every reserve assignee sits on an edge of $G(e)$, hence $\sigma(i)+k\ge e-\delta$ and $e_t^+(\mu)\ge e_o(\mu)-\delta$ for nonempty $t$; empty categories are vacuous. Unlike A-STDA, the loop exits only with $|\mu_t|=\hat q^t$ for every $t$, so no surviving quota is left standing empty.

\emph{Non-wastefulness.} At termination $|\mu_t|=\hat q^t$ for every $t$, since otherwise another de-reservation would have occurred, and $\mu_o$ holds the top $\min\{\hat q^o,|I|\}$ students, so either everyone is matched or $o$ is full. No unmatched student can be added anywhere within $\hat q$.
\end{proof}

\begin{proposition}[PA-ODR is fair to unmatched claimants]
\label{prop:paodr-fair}
The output of PA-ODR satisfies Definition~\ref{def:fair} restricted to unmatched claimants: there are no $t\in T$ and students $i,j$ with $t\in\T(i)\cap\T(j)$, $\sigma_t(i)>\sigma_t(j)$, $j\in\mu_t$ and $i\notin M(\mu)$.
\end{proposition}

\begin{proof}
Fix the final iteration, with cutoff $e$ and output $\mu$, and suppose $t\in T$ and $i,j$ satisfy $t\in\T(i)\cap\T(j)$, $\sigma_t(i)>\sigma_t(j)$, $j\in\mu_t$ and $i\notin M(\mu)$. First, $(i,t)$ is an edge of $G(e)$: being unmatched, $i$ is on the left side; $(j,t)$ is an edge, so $\sigma(j)+k\ge e-\delta$, and $\sigma_t(i)>\sigma_t(j)$ gives $\sigma(i)+k\ge e-\delta$. This is the crucial step: a student outranking a matched student inside a category is automatically admissible at the prevailing cutoff, so the cutoff never shields a matched student from a stronger claimant. Replacing $j$ by $i$ in $\mu_t$ now gives a capacitated matching on $G(e)$ of the same cardinality whose sorted score vector is lexicographically strictly larger, contradicting criteria 1 and 2.
\end{proof}

No capacity qualifier attaches to this proposition, unlike Proposition~\ref{prop:paodr-basic}. The step that carries it is the first one above: a student outranking a matched student inside a category is automatically admissible at the prevailing cutoff, so the cutoff never shields a matched student from a stronger claimant, and replacing one by the other cannot break score-feasibility either. Unlike SPSE the guarantee comes from optimality inside the final feasibility graph rather than from rejection monotonicity, and by Proposition~\ref{prop:paodr-unfair} the restriction to unmatched claimants cannot be lifted.

\subsection{What PA-ODR Gives Up}
\label{sec:paodr-limits}

Enforcing score-feasibility is not free; the price falls on fairness and incentives, and one source explains both: PA-ODR settles the reserve assignment by cardinality and score first, consults preferences only as a tie-breaker, and fills the open category without consulting them at all.

\begin{example}[PA-ODR is unfair to a matched claimant]
\label{ex:paodr-envy}
Let $T=\{t_1\}$, $q^o=q^{t_1}=1$, $k=5$, $\delta=10$, with $i_1(100,\{t_1\})$ and $i_2(90,\{t_1\})$ both ranking $t_1$ above $o$. PA-ODR gives $o$ the top raw score, so $\mu_o=\{i_1\}$ and $e=100$; as $95\ge90$, $(i_2,t_1)$ is an edge, $\mu_{t_1}=\{i_2\}$, every capacity is met and the algorithm stops without de-reserving, score-feasibly ($e_o-e_{t_1}^+=5\le\delta$). Yet $i_1$ sits in $o$ while strictly preferring $t_1$, whose occupant has $\sigma_{t_1}(i_2)=95<105=\sigma_{t_1}(i_1)$. The triple $(i_1,t_1,i_2)$ blocks Definition~\ref{def:fair} through clause (iv), and only through it: every student is matched, so the restriction of Proposition~\ref{prop:paodr-fair} is silent here.

The score rule does not excuse the failure, which matters because those notions were left unweakened. Honouring the claim gives $\mu_{t_1}=\{i_1\}$, $\mu_o=\{i_2\}$ with $e_o=90$ and $e_{t_1}^+=105$, score-feasible, so the complaint survives the score-relative reading. It is also what SPSE returns here, both students applying to $t_1$, which retains $i_1$. The mechanisms differ purely in whose claim they honour.
\end{example}

\begin{proposition}[PA-ODR is unfair to matched claimants]
\label{prop:paodr-unfair}
The output of PA-ODR need not be fair within categories in the sense of Definition~\ref{def:fair}. This already happens with $|T|=1$, and it happens with the complaint confined to the reserve categories, in an instance where every quota is filled and nothing is de-reserved. Nor need it be across-category efficient (Definition~\ref{def:acrosseff}).
\end{proposition}

\begin{proof}
The first claim is Example~\ref{ex:paodr-envy}. For the second let $T=\{t_1,t_2\}$, $q^o=q^{t_1}=q^{t_2}=1$, $k=5$, $\delta=10$, with $i_1(100,\emptyset)$, $i_2(95,\{t_1,t_2\})$, $i_3(90,\{t_2\})$ and $t_2\succ_{i_2}t_1\succ_{i_2}o$. Then $\mu_o=\{i_1\}$, $e=100$, admissibility needs $\sigma(i)\ge85$, and $G(100)$ carries $(i_2,t_1)$, $(i_2,t_2)$, $(i_3,t_2)$. Only $\{(i_2,t_1),(i_3,t_2)\}$ has cardinality $2$, so criterion 1 alone fixes the outcome, every quota is met and the loop exits at $\hat q=q$. But $\sigma_{t_2}(i_2)=100>95=\sigma_{t_2}(i_3)$ with $i_3\in\mu_{t_2}$ and $t_2\succ_{i_2}t_1=\mu(i_2)$: a within-blocking triple in which the open category plays no part. Honouring it by $\mu_{t_2}=\{i_2\}$, $\mu_{t_1}=\emptyset$ gives $e_o=100=e_{t_2}^+$, so this complaint too survives the score-relative reading.

For the last claim, take Example~\ref{ex:paodr-envy} and let $i_2$ rank $o$ above $t_1$ instead. Nothing in the run changes, PA-ODR consulting no preference before the reserve tie-break, so again $\mu_o=\{i_1\}$ and $\mu_{t_1}=\{i_2\}$; but now $t_1=\mu(i_2)\succ_{i_1}o=\mu(i_1)$ and $o=\mu(i_1)\succ_{i_2}t_1$, and the exchange gives $e_o=90$, $e_{t_1}^+=105$, hence a feasible and score-feasible matching. The pair $(i_1,i_2)$ blocks by exchange.
\end{proof}

The failure is structural, not a defect of the tie-breaker: criterion 1 maximizes reserve placements, and a strong overlapping student is the cheapest cover for a category only she can reach, so she is spent where she is scarce rather than where she is wanted. The last two clauses fail independently: one is a priority claim inside $T$, the other an exchange across $o$, which no criterion sees because $o$ is filled before the graph is built. What criterion 3 does secure is Definition~\ref{def:acrosseff} \emph{restricted to $T$}. Were $i\in\mu_c$ and $j\in\mu_{c'}$ with $c,c'\in T$ an exchange-blocking pair, admissibility at $e$ being category-independent makes $(i,c')$ and $(j,c)$ edges too, so exchanging them keeps the matched set and leaves criteria 1 and 2 indifferent while the sorted rank vector falls in two coordinates and rises in none, contradicting criterion 3.

\begin{proposition}[PA-ODR is not strategy-proof]
\label{prop:paodr-notsp}
There is an instance in which a student obtains a strictly $\succ_i$-better assignment under PA-ODR by misreporting her category preference.
\end{proposition}

\begin{proof}
Let $T=\{t_1,t_2\}$, $q^o=q^{t_1}=q^{t_2}=1$, $k=5$, $\delta=10$, with $i_1(100,\emptyset)$, $i_2(95,\{t_1,t_2\})$, $i_3(94,\{t_1,t_2\})$ and true preferences $t_1\succ_{i_2}t_2\succ_{i_2}o$, $o\succ_{i_3}t_1\succ_{i_3}t_2$. In every run $\mu_o=\{i_1\}$, $e=100$, and both $i_2$ and $i_3$ clear the bar, so $G(100)$ has all four edges between $\{i_2,i_3\}$ and $\{t_1,t_2\}$ and by Lemma~\ref{lem:matroid} both are matched. Criterion 2 is silent and criterion 3 decides between $\alpha:i_2\mapsto t_1,\ i_3\mapsto t_2$ and $\beta:i_2\mapsto t_2,\ i_3\mapsto t_1$. Truthfully the sorted rank vectors are $(3,1)$ for $\alpha$ and $(2,2)$ for $\beta$, so criterion 3 picks $\beta$ and $i_2$ gets $t_2$. Let $i_2$ instead report $t_1\succ o\succ t_2$, demoting $t_2$ below $o$ while keeping $t_1$ on top: the vectors become $(3,1)$ and $(3,2)$, which agree in the first coordinate, so criterion 3 picks $\alpha$ and $i_2$ gets $t_1$. As truly $t_1\succ_{i_2}t_2$, the misreport makes her strictly better off.
\end{proof}

This is the manipulation any egalitarian tie-breaker invites: reporting a category as worse than it is makes every assignment giving it to her look worse, which under a min--max rule pays whenever she would otherwise be sacrificed. It is no artifact of criterion 3 but reflects that PA-ODR is not a deferred-acceptance mechanism.

% \section{An Additional PA-ODR Example}

% It is useful to complement the previous example, which required de-reservation, with a case in which the mechanism stops immediately.

% \begin{example}[No de-reservation needed]
% Let
% \[
% T=\{t_1,t_2\},\qquad q^o=q^{t_1}=q^{t_2}=1,\qquad \delta=10,
% \]
% and let the common reserve bonus be $k=5$.
% Suppose the students are
% \[
% i_1(100,\emptyset),\quad i_2(88,\{t_1,t_2\}),\quad i_3(84,\{t_1\}),\quad i_4(86,\{t_2\}).
% \]
% Student $i_2$ prefers $t_2$ to $t_1$ and both reserve categories to the open category. Student $i_3$ prefers $t_1$ to $o$, and student $i_4$ prefers $t_2$ to $o$.

% In the first iteration, the open category takes $i_1$, so the open cutoff is $e=100$. Reserve feasibility therefore requires adjusted score at least $90$. Student $i_2$ satisfies
% \[
% 88+5=93\ge 90,
% \]
% so both edges $i_2\to t_1$ and $i_2\to t_2$ are feasible. Student $i_3$ does not satisfy the threshold for $t_1$ because
% \[
% 84+5=89<90,
% \]
% while student $i_4$ does satisfy the threshold for $t_2$ because
% \[
% 86+5=91\ge 90.
% \]

% Hence the feasible reserve graph contains edges
% \[
% i_2\to t_1,\qquad i_2\to t_2,\qquad i_4\to t_2.
% \]
% The maximum-cardinality reserve assignment has size $2$ and is unique:
% \[
% i_2\to t_1,\qquad i_4\to t_2.
% \]
% All reserve capacities are filled, so the algorithm stops immediately with
% \[
% \mu_o=\{i_1\},\qquad \mu_{t_1}=\{i_2\},\qquad \mu_{t_2}=\{i_4\}.
% \]

% The final cutoffs are
% \[
% e_o(\mu)=100,\qquad e_{t_1}^+(\mu)=93,\qquad e_{t_2}^+(\mu)=91.
% \]
% Thus
% \[
% e_o(\mu)-e_{t_1}^+(\mu)=7\le 10,\qquad e_o(\mu)-e_{t_2}^+(\mu)=9\le 10.
% \]
% So the outcome is score-feasible and no de-reservation occurs.
% \end{example}

% This example clarifies that de-reservation is not a built-in bias toward shrinking reserve quotas. It occurs only when some reserve capacity cannot be filled while preserving the cutoff rule.

\section{\uppercase{Prototype}}
\label{sec:experiments}

We coded the three mechanisms in Python and ran them on two families of $100$ randomized instances, seed $100$. Raw scores are uniform on $[0,100]$, the common bonus on $[0,15]$, $\delta=8$ unless stated, preferences over $V_i$ uniformly random, and each category enters $\T(i)$ independently with probability $0.45$ (a student drawing none is given one at random). Capacities are uniform on $q^o\le q^o_{\max}$, $q^t\le q^t_{\max}$, with $(|I|,|T|,q^o_{\max},q^t_{\max})=(200,6,30,15)$ and $(60,6,30,15)$, offering on average $47.99$ and $47.16$ reserve seats. The second family is the point: in the first, capacity falls short of $|I|$ and the mechanisms cannot separate on coverage; the second inverts this, average capacity $64.25$ against $|I|=60$. Every property is checked against its definition.

\begin{table*}[t]
\centering
\caption{Prototype outcomes over $100$ randomized instances per family. Omitted because they hold in all $600$ runs: feasibility and non-wastefulness (Def.~\ref{def:nw})---A-STDA's failure of the latter needs the configuration of Proposition~\ref{prop:tada-nw} and no random draw produced it. Percentages are of the family's average reserve seats.}
\label{tab:prototype}
\footnotesize
\setlength{\tabcolsep}{3pt}
\begin{tabular}{|L{0.14\textwidth}|L{0.08\textwidth}|r|r|r|r|r|L{0.13\textwidth}|}
\hline
Family & Mech. & Med.\ (ms) & Score-feas. & Within-fair & Acr.-eff. & Matched & De-reserved \\
\hline
$|I|=200$, $|T|=6$ & SPSE & 1.06 & 54/100 & 100/100 & 100/100 & 62.97 & 0.00 \\
(capacity-poor) & A-STDA & 1.66 & 100/100 & 100/100 & 100/100 & 58.18 & 0.00 \\
 & PA-ODR & 23.79 & 100/100 & 0/100 & 1/100 & 62.97 & 18.57 (39\%) \\
\hline
$|I|=60$, $|T|=6$ & SPSE & 0.14 & 45/100 & 100/100 & 100/100 & 55.48 & 0.00 \\
(capacity-rich) & A-STDA & 0.30 & 100/100 & 100/100 & 100/100 & 46.83 & 0.00 \\
 & PA-ODR & 8.75 & 100/100 & 63/100 & 63/100 & 56.61 & 44.09 (93\%) \\
\hline
\end{tabular}
\end{table*}

This prices the trade-off Proposition~\ref{prop:impossibility} forces rather than measuring average-case performance. First, SPSE breaches the cutoff rule in $46$ and $55$ of the two families' instances, so Proposition~\ref{prop:spse-not-sf} is no artifact of three-student examples.

Second, A-STDA's waste is real but does not show up as a violated axiom. It satisfies Definition~\ref{def:nw} in all $200$ instances, Proposition~\ref{prop:tada-nw} needing a configuration no random draw produced; what it does show is a shortfall in coverage, $58.18$ students against SPSE's $62.97$ on the capacity-poor family and $46.83$ against $55.48$ on the capacity-rich one. It leaves seats standing empty in $46$ and $54$ instances, roughly when SPSE would have used them illegally.

Third, fairness is where the mechanisms part, and the two notions locate PA-ODR's losses differently. SPSE and A-STDA keep both in all $600$ runs. PA-ODR is exchange-blocked in $99$ and $37$ instances, and in every one of them the blocking pair straddles the open category---no pair of reserve categories is ever exchange-blocked, the guarantee of Section~\ref{sec:paodr-limits} and exactly its boundary. Within categories it fails on priority instead, in $100$ and $37$ instances, so neither side of the design escapes: the open-first rule costs the exchanges and the cardinality-first rule costs the priorities.

Fourth, the two responses to Proposition~\ref{prop:impossibility} charge in different currencies. Where capacity binds, PA-ODR matches as many students as SPSE and pays purely in quota, $39\%$ of the reserve seats, while A-STDA pays purely in coverage, $4.79$ students, and keeps every seat. Where capacity is slack the gap widens: PA-ODR matches $56.61$ students, more than SPSE's $55.48$ and exactly the $\min\{\sum_c q^c,|I|\}$ bound it saturates by construction---but it costs $93\%$ of the reserve seats to buy $1.13$ students, while A-STDA retires nothing and matches $46.83$. Whether that is a good bargain is a policy question; the model's contribution is that it must be struck.

Fifth, the price tracks $\delta$: on the capacity-poor family, $\delta=2,4,8,15,25$ leave SPSE score-feasible in $34,43,54,79,97$ instances and retire $30.32,26.38,18.57,7.84,0.61$ reserve seats, so at a permissive gap all three nearly coincide.

Two caveats. Running times are medians on one laptop, reported only to show that nothing here is intractable; A-STDA terminated in two rounds in every instance. And the breach rates mix the two causes of Section~\ref{sec:spse}: lowering the overlap probability to $0.15$ leaves SPSE score-feasible in only $6$ of $100$ capacity-poor instances, worse rather than better, thin eligibility stranding weak candidates. The breach rate is therefore not a measure of the overlap-specific cause of Proposition~\ref{prop:spse-not-sf}.

\section{\uppercase{Conclusion}}

Overlapping reserve categories create a genuine allocation problem because admitted students count toward exactly one category. Score constraints make feasibility depend on endogenous cutoffs, so non-wastefulness must be read relative to the constraint, and full capacity utilization can fail even when every quota is individually honourable. Proposition~\ref{prop:impossibility} forces a choice: score-feasibility, open-completeness and reserve-completeness cannot be had together. SPSE gives up score-feasibility, keeping the deferred-acceptance package; A-STDA gives up reserve-completeness, and with it seats and incentives (Propositions~\ref{prop:tada-nw}, \ref{prop:tada-notsp}); PA-ODR gives up the announced capacities. The two fairness notions divide the three differently from the score rule: SPSE and A-STDA keep both, PA-ODR neither, and it fails them for different reasons. Maximizing reserve placements spends a strong overlapping student where she is scarce rather than where she is wanted (Proposition~\ref{prop:paodr-unfair}), while filling the open category first, before any preference is consulted, leaves every blocking exchange straddling it. Tables~\ref{tab:comparison} and~\ref{tab:prototype} record the rest.

We should not overstate what overlap explains: cutoff-gap breaches and order dependence occur without it. The specifically overlapping content is the routing of multiply eligible students and the Hall obstruction of Proposition~\ref{prop:impossibility}: the gap between Hall's condition on singleton category sets and on larger ones.

% \subsection{Alternative Fairness Formulations}
% It is useful to record several nearby fairness formulations, because they reveal what is conceptually difficult in the model.

% \subsubsection{Open-category-inclusive fairness}
% One concern is that a fairness definition restricted only to reserve categories may miss objections involving the open category. If the open category is not explicitly included, a high-score student who should obviously displace a low-score open assignee may fail to generate a violation under a narrowly stated reserve-only notion.

% For this reason, consider enlarging each student's relevant category set to
% \[
% \T^+(i):=\T(i)\cup\{o\}.
% \]
% Under this convention, within-category fairness can be imposed uniformly across all seat categories, including the open category. The corresponding idea is:

% \begin{quote}
% if a lower-score student occupies some category $c\in T\cup\{o\}$, then no higher-score student who is eligible for $c$ should remain unmatched.
% \end{quote}

% This formulation is stronger than reserve-only within-category fairness and is particularly natural when the open category is viewed as just another counting category in the one-to-one framework.

% \subsubsection{Cardinality-based non-wastefulness}
% Another variant defines non-wastefulness through the impossibility of increasing the number of matched students while preserving already matched students and maintaining score-feasibility under the final de-reserved capacity vector. Conceptually, this version says that an outcome is non-wasteful if there exists no feasible reconfiguration that keeps all current winners and adds somebody else on top.

% This cardinality-based definition is well aligned with de-reservation mechanisms because the relevant capacity vector changes over time. It avoids describing wastefulness purely in terms of the original capacity vector, which may no longer be the correct object after some reserve categories have been transferred to the open category.

% \subsection{Implementation and Empirical Questions}
% % Even without a full empirical section, it is useful to state what future work would study.
% There are several dimensions worth investigation when considering empirical implementation. The first natural dimension is the degree of overlap. If most students qualify for at most one reserve category, then the overlapping problem is close to a standard reserve system. If overlap is common, then the mechanism's counting rule becomes much more important.

% The second dimension is the magnitude of the common reserve bonus. When the bonus is large, the feasible edge set under adjusted-cutoff score-feasibility expands; when it is small, de-reservation occurs more often.

% The third dimension is the strictness of the score-gap parameter $\delta$. For example, in PA-ODR, when $\delta$ is large, many more reserve assignments are feasible and de-reservation will occur less often. When $\delta$ is small, the mechanism may repeatedly transfer reserve capacity into the open category.

% Finally, one could compare output statistics across mechanisms: number of matched students, number of categories de-reserved, overlap usage patterns, and frequency of fairness complaints under ex post diagnostics. The structure developed in this paper provides the language for such an exercise though the actual simulations have not yet been carried out.

% \section{Mechanism Comparison}

% The analysis ultimately points toward a comparison between two broad approaches.

% \begin{center}
% \renewcommand{\arraystretch}{1.3}
% \begin{tabular}{lll}
% \toprule
% \textbf{Property} & \textbf{SPSE} & \textbf{PA-ODR} \\
% \midrule
% One-to-one counting & Yes & Yes \\
% Feasibility & Yes & Yes \\
% Non-wastefulness & Yes & Yes, w.r.t.\ final $\hat q$ \\
% Within-category fairness & Yes & Intended under final feasibility graph \\
% Replacement fairness & Yes & Intended under final feasibility graph \\
% Score-feasibility & Not guaranteed & Enforced by construction \\
% Capacity reallocation & Fixed at $q$ & Endogenous \\
% Uses category preferences & Through proposals & Through final lexicographic tie-breaking \\
% \bottomrule
% \end{tabular}
% \end{center}

% This comparison captures the main design contrast. SPSE keeps the initial reserve capacities fixed and has clean deferred-acceptance style fairness properties, but it does not generally satisfy the score-gap rule. PA-ODR enforces the score-gap rule by allowing de-reservation and therefore endogenizing the effective reserve capacities.

Several questions remain. The sharpest is whether the incentive loss is forced: both mechanisms enforcing the cutoff rule are manipulable (Propositions~\ref{prop:tada-notsp}, \ref{prop:paodr-notsp}) and the strategy-proof one abandons it, so an incentive counterpart to Proposition~\ref{prop:impossibility} may be within reach. Second, whether minimum de-reservation is computable efficiently for unbounded $|T|$, where we conjecture hardness, realizability resting on a Hall-type condition over all subsets of $T$; Cs\'aji and Sun \cite{CsSu26a} solve the nearest fixed-market problem in polynomial time, while overlapping types with per-type quotas are $\Sigma_2^{\mathrm{P}}$-complete \cite{CGH20a}. Third, whether routing is essential at all: one could apply the bonus and admit purely by adjusted score. Finally, the common bonus is deliberate: category-specific bonuses would break the property that an overlapping student is admissible for all her categories or none.

% \section{Conclusion}

% The underlying message is consistent throughout. Overlapping reserve categories create a genuine allocation problem because admitted students must be counted in exactly one category. Score constraints complicate the problem further by making feasibility depend on endogenous cutoffs. As a result, ordinary notions of wastefulness and fairness must be redefined to incorporate the score rule, and stronger merit requirements can even make existence fail.

% The analysis then develops two mechanisms. SPSE combines common reserve-side score elevation with deferred-acceptance logic and keeps the initial reserve capacities fixed, but it need not satisfy score-feasibility. PA-ODR instead combines common reserve-side score elevation with an open-first rule and converts unfulfillable reserve capacity into open capacity when necessary. The first uses a fixed capacity vector; the second enforces the score constraint through endogenous capacity reallocation.

% The unresolved tension between fixed capacities and score-feasibility is therefore not an accidental artifact of a particular algorithm. It appears to be a core design challenge of admissions with overlapping reserve categories under cutoff restrictions.

\bibliographystyle{apalike}
\bibliography{references}

\end{document}
